% Figure: thermal buckling of an axially restrained bar. A slender bar held
% between two rigid walls is heated. It would expand freely by alpha*L*Delta T,
% but the walls forbid the expansion and instead build up a compressive axial
% force. When that force reaches the Euler buckling load the straight bar gives
% way to a bowed shape. The plot shows the compressive force rising linearly
% with temperature until it meets the flat Euler ceiling, fixing the critical
% temperature rise.
\begin{figure}[htbp]
\centering
\begin{tikzpicture}[
  wall/.style={draw=engblack, thick},
  bar/.style={draw=engblue, very thick},
  axis/.style={draw=engblack, -{Stealth}, thick},
  force/.style={draw=engred, very thick},
  ceil/.style={draw=engamber, dashed, thick},
  guide/.style={draw=enggray, dashed, thin},
  label/.style={font=\small},
]
% ===== restrained bar sketch (top) =====
\draw[wall] (0.2,2.4) -- (0.2,3.4);
\foreach \y in {2.5,2.7,...,3.3} {\draw[enggray, thin] (0.2,\y) -- (0.0,\y-0.12);}
\draw[wall] (5.8,2.4) -- (5.8,3.4);
\foreach \y in {2.5,2.7,...,3.3} {\draw[enggray, thin] (5.8,\y) -- (6.0,\y-0.12);}
% straight bar before buckling (faint)
\draw[enggray, thin] (0.2,2.9) -- (5.8,2.9);
% buckled bowed bar
\draw[bar] plot[domain=0.2:5.8, samples=80]
  ({\x}, {2.9 + 0.45*sin( deg( (\x-0.2)/5.6*3.14159 ) )});
\node[engblue, font=\scriptsize, anchor=south] at (3.0,3.45)
  {buckled shape, $\Delta T>\Delta T_\mathrm{cr}$};
% arrows showing thermal push against walls
\draw[-{Stealth}, engred, thick] (0.55,2.9) -- (0.05,2.9);
\draw[-{Stealth}, engred, thick] (5.45,2.9) -- (5.95,2.9);
\node[engred, font=\scriptsize] at (3.0,2.55) {restrained expansion $\Rightarrow$ compression $N$};

% ===== force vs temperature plot (bottom) =====
\begin{scope}[yshift=-0.2cm]
\draw[axis] (0,-2.3) -- (0,0.9) node[anchor=south east, label] {axial force $N$};
\draw[axis] (0,-2.3) -- (6.4,-2.3) node[anchor=west, label] {temperature rise $\Delta T$};
% N = EA alpha Delta T, linear rising line from origin (at bottom-left)
% origin of plot at (0,-2.3). Slope chosen so line meets ceiling at x=4.0.
\draw[force] (0,-2.3) -- (4.0,0.5);
\node[engred, font=\scriptsize, anchor=west, rotate=35] at (1.6,-1.35)
  {$N=EA\,\alpha\,\Delta T$};
% Euler ceiling: horizontal dashed line at N = P_cr (plotted y = 0.5)
\draw[ceil] (0,0.5) -- (6.2,0.5);
\node[engamber, font=\scriptsize, anchor=west] at (4.3,0.5)
  {$P_\mathrm{cr}=\pi^2EI/L^2$};
% critical temperature guide
\draw[guide] (4.0,-2.3) -- (4.0,0.5);
\fill[engred] (4.0,0.5) circle (2pt);
\node[font=\scriptsize, anchor=north] at (4.0,-2.35) {$\Delta T_\mathrm{cr}$};
\end{scope}
\end{tikzpicture}
\caption[Thermal buckling of a restrained bar]{A slender bar held between rigid
walls cannot expand when heated, so the forbidden thermal strain
$\alpha\,\Delta T$ is turned into an elastic compressive strain and the bar
develops an axial force $N=EA\,\alpha\,\Delta T$ that rises in direct
proportion to the temperature rise. When $N$ reaches the Euler buckling load
$P_\mathrm{cr}=\pi^2EI/L^2$ the straight bar loses stability and bows out. The
critical temperature rise $\Delta T_\mathrm{cr}=\pi^2 I/(A\alpha L^2)$ where the
rising force line meets the flat Euler ceiling depends on the slenderness of
the bar, not on the modulus, which cancels. The energy method reads the
buckling condition off the second variation of the total potential energy, the
same criterion the mechanical column-buckling example used.}
\label{fig:vol03ch05:thermal-buckling}
\end{figure}
